\section{Fresh inventory blocks and Brownian contacts}
\label{arm:inputs}

We first establish the contour estimates used in local graph reconstruction
and in the later control of discrete endpoints.  A fresh inventory block
has the same Brownian limit uniformly over its incoming stack.  Brownian
survival estimates then rule out specific contact configurations, including
contacts whose endpoints are chosen after the contour is observed.
All parameters are fixed as in Equation~\eqref{bg:parameters}; no passage
metric or graph-distance tightness is used in this section.

\subsection{Bounds determined by a fresh raw word}
\label{arm:raw-word-bounds}

Write the inventory alphabet as $\{h,c,H,C,F\}$, with $h,c$
producing the two burger types, $H,C$ ordering the corresponding type,
and $F$ ordering the freshest available burger.  The symbol probabilities
are $1/4,1/4,(1-b_F)/4,(1-b_F)/4,b_F/2$, respectively.
For a finite word $w=w_1\cdots w_m$, let $R(w)$ be its reduction
from an empty burger stack, retaining unsatisfied orders permanently.
Set
\[
 L_k(w)=|R(w_1\cdots w_k)|,\qquad
 M_m(w)=\max_{0\le k\le m}L_k(w),\qquad
 F_k(w)=N_F(R(w_1\cdots w_k)).
\]
An \emph{exterior typing} $\sigma$ assigns either type to every $F$
which this procedure leaves unmatched.  Every other $F$ retains the type
of its internally matched burger.  Let $Z^\sigma(0)=0$ and give each
burger the increment $e_i$ of its type and each resolved order the
increment $-e_i$.  We interpolate linearly between integer times.
The exterior typing may depend on the whole word and on auxiliary data.
Every genuine initial past induces an exterior typing.

\begin{lemma}[Deterministic reduction identities]
\label{arm:raw-identities}
For every finite word and every exterior typing,
\begin{equation}
\label{arm:raw-length-identity}
 L_k(w)=\sum_{i=1}^2\left(Z_i^\sigma(k)
             -2\min_{0\le j\le k}Z_i^\sigma(j)\right).
\end{equation}
Consequently,
\begin{align}
 M_m(w)&\le 2\sum_{i=1}^2
       \left(\max_{0\le j\le m}Z_i^\sigma(j)
                    -\min_{0\le j\le m}Z_i^\sigma(j)\right),
       \label{arm:raw-contour-oscillation}\\
 M_m(w)&\le 6\max_{i\in\{1,2\},\,0\le j\le m}|Z_i^\sigma(j)|,
       \label{arm:raw-contour-tightness}\\
 \max_{0\le j,k\le m}\|Z^\sigma(k)-Z^\sigma(j)\|_1&\le 2M_m(w).
       \label{arm:raw-universal-domination}
\end{align}
For finite words $u,v$,
\begin{equation}
\label{arm:raw-subadditivity}
 |R(uv)|\le |R(u)|+|R(v)|.
\end{equation}
Finally, $F_k(w)$ is nondecreasing in $k$, and for any two exterior
typings $\sigma,\tau$,
\begin{equation}
\label{arm:raw-typing-error}
 \max_{0\le k\le m}\|Z^\sigma(k)-Z^\tau(k)\|_1\le 2F_m(w).
\end{equation}
\end{lemma}

\begin{proof}
Every burger produced inside the word is fresher than every burger in
an initial past.  Induction over symbols therefore shows that the
available internal burgers and their internal matching are independent
of that past.  A typed order uses an internal burger of its type whenever
one is available.  An $F$ uses the freshest internal burger whenever the
internal stack is nonempty.  A step which instead uses the past does not
change the internal stack.  In particular, each externally unmatched
$F$ occurs when the entire internal burger stack is empty.  Replacing
only these $F$'s by arbitrary typed orders therefore preserves every
internal pair and every unpaired position.

In the resulting fully typed word let $O_i(k)$ and $B_i(k)$ count the
unpaired orders and burgers of type $i$.  An order is unpaired exactly
when the corresponding coordinate reaches a new negative running
minimum.  Thus
\[
 O_i(k)=-\min_{0\le j\le k}Z_i^\sigma(j),\qquad
 B_i(k)=Z_i^\sigma(k)-\min_{0\le j\le k}Z_i^\sigma(j).
\]
Summing $O_i+B_i$ proves
Equation~\eqref{arm:raw-length-identity}.
Since $Z^\sigma(0)=0$, its summands are bounded by twice the
corresponding coordinate oscillations and by three times the maximum
absolute coordinate height.  This proves
Equations~\eqref{arm:raw-contour-oscillation}
and~\eqref{arm:raw-contour-tightness}.
Also $Z_i^\sigma=B_i-O_i$, so
$\|Z^\sigma(k)\|_1\le L_k(w)$; the triangle inequality proves
Equation~\eqref{arm:raw-universal-domination}.

For Equation~\eqref{arm:raw-subadditivity}, process $v$ on the retained
burger stack of $u$.  All its internal pairs persist.  Its orders which
were internally unpaired either remain unpaired or consume a retained
burger of $u$.  Such a consumption removes two letters from the two
separate reductions; no new letter is created.

An order unpaired in a prefix stays unpaired in every longer prefix,
since orders only consume earlier burgers.  Hence $F_k$ is
nondecreasing.  The increments for two exterior typings can differ
only at these unpaired $F$ positions, and each difference has
$\ell^1$ norm two.  Summing the increment differences gives
Equation~\eqref{arm:raw-typing-error}.
\end{proof}

\begin{proposition}[Exponential bound from square-root tightness]
\label{arm:raw-exponential}
For the iid raw word with fixed $b_F\in(0,1/2)$ there are constants
$c,C>0$ such that, for every integer $m\ge1$ and $x\ge0$,
\begin{equation}
\label{arm:raw-exponential-tail}
 \mathbb P[M_m>x\sqrt m]\le Ce^{-cx}.
\end{equation}
In particular, for some $\lambda>0$ and $K<\infty$,
\begin{equation}
\label{arm:raw-exponential-moment}
 \sup_{m\ge1}\mathbb E\exp\{\lambda M_m/\sqrt m\}\le K.
\end{equation}
The same bounds hold conditionally on any previously observed
information under which the raw block remains iid with this law.
\end{proposition}

\begin{proof}
Adjoin an independent bilateral iid past and resolve the block using
that past.  Sheffield's invariance principle
\cite[Theorem~2.5]{SheffieldInventory}, equivalently
\cite[Theorem~1.5]{GwynneMaoSunConeTimes}, gives tightness of the family
\[
 m^{-1/2}\max_{i,j\le m}|Z_i(j)|.
\]
Equation~\eqref{arm:raw-contour-tightness} gives tightness of the
block-only random variables $M_m/\sqrt m$.
Choose $A>0$ such that for every $m\ge1$, with
$a_m=\lceil A\sqrt m\rceil$,
\[
 \mathbb P[M_m\ge a_m]\le\tfrac12.
\]
Fix $m$, write $a=a_m$, and, in an infinite iid raw word, restart
whenever the current block first has reduced length $a$.
More precisely, let $S_0=0$ and
\[
 S_{r+1}=\inf\{t>S_r:|R(X_{S_r+1}\cdots X_t)|=a\}.
\]
A single added symbol changes reduced length by one in absolute
value, so these stopping times have no overshoot.  They are finite
almost surely: reduced length dominates absolute net burger count,
which is a simple symmetric random walk.  Strong Markov for the raw
iid symbols makes $S_r-S_{r-1}$ iid.  Put
$N_m=\max\{r:S_r\le m\}$.  By
Equation~\eqref{arm:raw-subadditivity}, every prefix up to time $m$
has reduced length at most $a(N_m+1)$.  Therefore, for integers $r\ge0$,
\[
 \mathbb P[M_m>ra]\le\mathbb P[N_m\ge r]
 \le\prod_{j=1}^r\mathbb P[S_j-S_{j-1}\le m]
 \le2^{-r}.
\]
Since $a_m/\sqrt m\le A+1$, take
$r=\lfloor x\sqrt m/a_m\rfloor$ to obtain
Equation~\eqref{arm:raw-exponential-tail} with
$C=2$ and $c=\log(2)/(A+1)$.  For $0<\lambda<c$, integration of the
tail gives $K\le1+2\lambda/(c-\lambda)$.
The estimate depends only on the raw block, proving its conditional
version.  Iterating this conditional moment estimate also bounds the
exponential moment of a sum over successively exposed fresh blocks
whose lengths are chosen before the respective blocks are inspected.
\end{proof}

Sheffield already observed that the invariance principle remains valid
for an arbitrary fixed initial burger stack
\cite[Remark~3.17]{SheffieldInventory}. The next statement uses the
quantitative unmatched-order estimate of
\cite[Lemma~3.7]{GwynneMaoSunConeTimes} to make the comparison uniform
over randomized, word-dependent exterior typings.

\begin{proposition}[Uniformity over exterior typings]
\label{arm:uniform-past-limit}
Fix $b_F\in(0,1/2)$.  Let $\nu_{b_F}$ be the law on
$C([0,1],\mathbb R^2)$ of Brownian motion started from zero with
covariance matrix at time $t$ equal to
\[
 \frac t2
 \begin{pmatrix}1-b_F&b_F\\b_F&1-b_F\end{pmatrix}.
\]
For the bounded-Lipschitz distance associated with the uniform
$\ell^1$ path metric,
\begin{equation}
\label{arm:uniform-past-bl}
 \sup_\sigma d_{\mathrm{BL}}\left(
   \mathcal L(m^{-1/2}Z^\sigma(m\,\cdot)),\nu_{b_F}\right)
   \longrightarrow0,
\end{equation}
where $\sigma$ ranges over all possibly randomized, word-dependent
exterior typings.  The convergence holds conditionally, uniformly
over previously observed data which leave the block fresh and iid.
A fixed finite family of deterministic disjoint blocks, whose lengths each
tend to infinity, converges jointly to independent Brownian paths after
applying the respective square-root normalizations.
\end{proposition}

\begin{proof}
The only additional public input is
\cite[Lemma~3.7, Equation~(3.15)]{GwynneMaoSunConeTimes}.
Put
\[
 \mu'=\frac{\pi}{2\left(\pi+
       \arctan\bigl(\sqrt{1-2b_F}/b_F\bigr)\right)}\in(1/3,1/2).
\]
That lemma says that for each $\nu>\mu'$,
$\mathbb P[\exists i\ge n:F_i\ge i^\nu]$ decays faster than every
negative power of $n$.  Choose $\nu\in(\mu',1/2)$.  It follows that
$F_m/\sqrt m\to0$ in probability.  Its monotonicity, already proved
in Lemma~\ref{arm:raw-identities}, simultaneously controls all prefixes.

Adjoin an independent bilateral past, and let $\tau$ be its genuine
typing.  Theorem~2.5 of \cite{SheffieldInventory} gives
$m^{-1/2}Z^\tau(m\,\cdot)\Rightarrow\nu_{b_F}$.
For each bounded $1$-Lipschitz test with absolute value at most one,
Equation~\eqref{arm:raw-typing-error} bounds the difference between
its expectations under $\sigma$ and $\tau$ by
\[
 \mathbb E\left[\min\{2,2F_m/\sqrt m\}\right]\longrightarrow0.
\]
This bound is simultaneous over all $\sigma$, proving
Equation~\eqref{arm:uniform-past-bl}.  Choose the auxiliary bilateral
past independently of the previously observed data to obtain the
conditional statement.  For deterministic disjoint blocks choose independent
auxiliary bilateral pasts for the respective blocks; their reference
contours are independent, and the same deterministic error bound
proves the finite-product assertion.
\end{proof}

\begin{remark}[Scope of fresh-block estimates]
\label{arm:fresh-block-scope}
Starting a block after a stopping time in the raw-symbol filtration
preserves the needed iid law.  Conditioning on a block's own stopping
duration, an exact terminal displacement, or a finite-word completion
event generally does not.  Such conditioning requires a separate
change-of-density argument.  The preceding statements only permit
arbitrary assignments to the raw-unmatched flexible orders; they do
not permit changing the type of an internally matched flexible order.
All constants and limits concern fixed $b_F<1/2$.
\end{remark}


The raw-word estimates control contour fluctuations without conditioning
on how the block is completed.  We next obtain the one- and two-wall
probability bounds needed for contacts and for the later address argument.

\subsection{Brownian survival and landing}

Let $Z=(Z_1,Z_2)$ be a centered Brownian motion with nondegenerate
covariance, equal coordinate variances, and correlation
$\rho=-\cos(\pi\gamma^2/4)$.  A fixed change of time normalizes the
variances.  Write
\[
 \theta=\frac{\pi\gamma^2}{4}\in(\pi/2,\pi),\qquad
 d_0=0,\quad d_1=\frac12,\quad d_2=\frac{\pi}{2\theta}.
\]
Whitening maps the positive quadrant to a wedge of opening $\theta$.
All constants in this subsection can depend on the fixed covariance.

\begin{lemma}[A wedge survival bound]\label{arm:wedge-survival}
Let $W$ be a planar wedge of opening $\theta\in(\pi/2,\pi)$, let
$\nu=\pi/\theta$, and let $\tau_W$ be the first exit time of standard
Brownian motion.  For $0<t_0\le t\le t_1<\infty$ there is a constant
$C$ such that
\[
 \mathbb P_z(\tau_W>t)\le C\min\{1,|z|^\nu\},\qquad z\in W.
\]
For one-dimensional Brownian motion killed at zero, the corresponding
bound is $C\min\{1,z\}$.
\end{lemma}

\begin{proof}
Brownian scaling reduces the assertion to $t=1$.  Put
$D=W\cap B(0,1)$ and use polar coordinates $(r,\phi)$, with
$0<\phi<\theta$.  The probability $q(r,\phi)$ of leaving $D$ through
its circular arc before a side has the harmonic expansion
\[
 q(r,\phi)=\frac4\pi
 \sum_{\substack{m\ge1\\m\text{ odd}}}
       \frac{r^{m\nu}}{m}\sin(m\nu\phi).
\]
For $r<1$ this is the harmonic extension of the sine series of the
constant function one on the circular arc, with zero side values.
One can first take partial sums on smaller concentric sectors and
then use uniqueness of bounded harmonic extensions.  The two corner
points have harmonic measure zero.  Absolute summation gives
$q(r,\phi)\le C r^\nu$ for $r\le1/2$; the bound for larger $r$ follows
from $q\le1$.

For the exit time of $D$, set
\[
 f(\phi)=\frac12\left(\frac{\cos(2\phi-\theta)}{\cos\theta}-1\right).
\]
Since $\cos\theta<0$ and
$\cos(2\phi-\theta)\ge\cos\theta$, one has $f\le0$.  The function
vanishes at the two sides and $f''+4f=-2$.  Choose $A<\infty$ such
that $-f(\phi)\le A\sin(\nu\phi)$; the ratio has finite limits at the
two endpoints because both numerator and denominator vanish linearly.
Then
\[
 v(r,\phi)=A r^\nu\sin(\nu\phi)+r^2 f(\phi)
\]
is nonnegative on $D$, satisfies $-\tfrac12\Delta v=1$, and obeys
$v\le A r^\nu$.  Nonnegativity follows from $\nu<2$ and
$r^2 f\ge r^\nu f$.  Applying the stopped It\^o formula on compact
subdomains, then increasing the subdomains to $D$, gives
$\mathbb E_z\tau_D\le v(z)\le A r^\nu$.
Therefore
\[
 \mathbb P_z(\tau_W>1)
 \le q(z)+\mathbb P_z(\tau_D>1)
 \le C r^\nu.
\]
The one-dimensional assertion follows from the reflection principle.
\end{proof}

\begin{lemma}[Uniform killed landing]\label{arm:killed-landing}
Constrain $k\in\{1,2\}$ coordinates of $Z$ to be nonnegative, and let
$\tau_k$ be the first violation.  For $0<h\le1$,
\begin{align}
 \sup_{z\in[0,h]^k}\mathbb P_z(\tau_k>1)
     &\le C h^{2d_k},\label{arm:survival-small}\\
 \sup_{z\in[0,\infty)^k}
 \mathbb P_z(\tau_k>1,\ Z(1)\in[0,h]^k)
     &\le C h^{k+2d_k}.\label{arm:landing-small}
\end{align}
For $k=1$, only the constrained coordinate is included in the starting
and landing vectors.  The same bounds hold on any fixed compact range
of positive durations, and conditionally when the starting position and
the constant lower walls are measurable before the Brownian increments
are sampled.
\end{lemma}

\begin{proof}
Whitening and Lemma~\ref{arm:wedge-survival} prove
Equation~\eqref{arm:survival-small}.  Let $p_t^D(x,y)$ be the killed
transition density in the half-line or whitened wedge.  Its symmetry,
the semigroup identity, and the bounded free Gaussian transition
density imply
\[
 p_1^D(x,y)
 =\int_Dp_{1/2}^D(x,z)p_{1/2}^D(z,y)\,dz
 \le C\int_Dp_{1/2}^D(z,y)\,dz
 =C\mathbb P_y(\tau_D>1/2).
\]
Integrating over the whitened landing set, which has volume $O(h^k)$
and lies at distance $O(h)$ from the corner, proves
Equation~\eqref{arm:landing-small}, uniformly in $x$.
Scaling handles the duration range.  Conditioning on past-measurable
starting data leaves the same Brownian increment law, proving the last
assertion.
\end{proof}

\begin{lemma}[Fresh lattice slots]\label{arm:lattice-slots}
Fix $\delta>0$.  For every sufficiently large fixed $R$, a fresh raw
inventory block of length $r(1+O(R^{-1}))$ satisfies the following
conditional bounds, uniformly over its raw past, once $r$ is sufficiently
large.  Put $h=R^{-1/2+\delta}$ and let the lower walls be measurable
before the block is exposed.
\begin{enumerate}[label=(\roman*)]
 \item Starting within $O(h\sqrt r)$ of $k$ walls and remaining above
       them throughout the block has probability at most
       $C R^{-(1-2\delta)d_k}$, with a further arbitrarily small
       exponent loss allowed.
 \item Starting anywhere above the walls, remaining above them, and
       landing within $O(h\sqrt r)$ of all $k$ walls has probability
       at most $C R^{-(1-2\delta)(k/2+d_k)}$, with the same convention.
\end{enumerate}
The constants in the $O(h\sqrt r)$ tolerances are fixed before $R$.
\end{lemma}

\begin{proof}
For fixed $R$, Proposition~\ref{arm:uniform-past-limit} gives uniform
weak convergence of the block increment paths.  Enlarge all path and
landing tolerances by a fixed small multiple of $h$; the corresponding
closed Brownian events bound the desired limsup.  Their probabilities
are bounded by Lemma~\ref{arm:killed-landing}, with a changed constant.
Uniformity over starting points in a fixed compact set follows by a
finite net of starting heights and one more tolerance enlargement.
For the landing test, starting points outside a sufficiently large
compact set require a correspondingly large block fluctuation;
Proposition~\ref{arm:raw-exponential} makes that probability smaller
than any prescribed fixed power of $R$ by choosing the compact set
after fixing $R$.  Starting below a wall has probability zero unless
an already allowed tolerance covers the discrepancy.  Finally choose
the block length large enough that the uniform weak-convergence error
is less than the desired small power of $R$.  This proves the stated
bounds, including conditioning on a fresh raw past.
\end{proof}


The following generic-contact statements also apply to the protected
charts.  Their proofs use only the Brownian estimates just established
and Gaussian increments between separated deterministic windows.

\subsection{Generic Brownian contacts}

Throughout this subsection, $Z=(L,R)$ is a two-sided Brownian motion with
unit coordinate variances and correlation $\rho\in(0,1)$.  A common constant
variance factor changes none of the assertions.  Set
\[
 \theta=\arccos(-\rho)\in(\pi/2,\pi),\qquad
 d_2=\frac{\pi}{2\theta}>\frac12.
\]
For a continuous scalar function $f$, a \emph{horizontal chord} is a pair
$s<t$ such that
\[
 f(s)=f(t)=\min_{[s,t]}f.
\]
A \emph{branch triple} consists of three distinct times $a<b<c$ for which
$f(a)=f(b)=f(c)=\min_{[a,c]}f$.  Thus its middle time is a local minimum.
All the following assertions concern nonzero chord spans and strictly
positive witnessing time intervals.

\begin{lemma}[Images of Brownian record times]\label{arm:record-height-images}
Fix a compact time interval $K$ and $\delta>0$.  The surrounding Brownian
path is retained far enough to include every arm of length $\delta$ used
below.  For signs $\sigma,\tau\in\{-1,1\}$, let
\[
 T_{\sigma,\tau}=
 \left\{t\in K:
 \begin{array}{l}
 L(t+\sigma u)\ge L(t),\\
 R(t+\tau u)\ge R(t)
 \end{array}
 \text{ for every }0\le u\le\delta\right\}.
\]
If $\sigma=\tau$, the upper box dimensions of $Z(T_{\sigma,\tau})$
and of either coordinate projection are at most $2(1-d_2)<1$.
If $\sigma\ne\tau$, these upper box dimensions are zero.
Furthermore, almost surely there is no time $t\in K$ at which $L$ has
lower-wall arms of length $\delta$ on both sides and $R$ has such an arm
on either one of the two sides.  The same statement holds with $L,R$
interchanged.
\end{lemma}

\begin{proof}
We first record the required survival bounds.  For fixed $v>0$ and
$0<r\le1$, the reflection principle gives
\[
 \mathbb P\left[\inf_{0\le u\le v}L(u)\ge-r\right]\le C_vr.
\]
Whitening the covariance sends the two lower walls to a wedge of opening
$\theta$.  Lemma~\ref{arm:wedge-survival} and Brownian scaling give
\begin{equation}\label{arm:generic-wedge-survival}
 \mathbb P\left[\inf_{0\le u\le v}L(u)\ge-r,
                  \inf_{0\le u\le v}R(u)\ge-r\right]
 \le C_v r^{2d_2}.
\end{equation}
Only this elementary Brownian survival estimate is used here; no lattice
or passage-kernel estimate is involved.

Partition $K$ into intervals $I=[s,s+h]$ of length $h=2^{-n}$, allowing
two extra intervals at its ends.  Fix $\varepsilon\in(0,1/2)$ and put
$r=h^{1/2-\varepsilon}$.  Almost surely, for all sufficiently large $n$,
both coordinate oscillations on every such interval are at most $r$.
For example this follows from Brownian H\"older continuity with exponent
$1/2-\varepsilon/2$.  Suppose such an interval contains a record time
$t$.  If a required arm points forward, then its constraint implies
\[
 \inf_{0\le u\le\delta/2}
       \{Z_i(s+h+u)-Z_i(s+h)\}\ge-r;
\]
if it points backward, the analogous constraint is
\[
 \inf_{0\le u\le\delta/2}
       \{Z_i(s-u)-Z_i(s)\}\ge-r.
\]
These implications hold when $h\le\delta/2$.  They use deterministic
walls $-r$ relative to the respective cell endpoints, so their probability
bounds do not condition on a selected record time.

For two arms on the same side, Equation
\eqref{arm:generic-wedge-survival} bounds this necessary exterior event
by $Cr^{2d_2}$.  For arms on opposite sides the two events use disjoint
time increments, and the reflection principle bounds their product by
$Cr^2$.  In the last assertion of the Lemma, the necessary exterior event
uses two coordinate arms on one side and one coordinate arm on the other;
its probability is at most $Cr^{2d_2+1}$.

Here is the resulting covering calculation, including the dimension
needed below.  If the survival power is $\alpha$, the number $N_h$ of
cells satisfying the necessary exterior event obeys
\begin{equation}\label{arm:record-cell-count}
 \mathbb E N_h\le C h^{-1}r^\alpha
          =C h^{-1+(1/2-\varepsilon)\alpha}.
\end{equation}
For $\alpha\le2$ and every $\zeta>0$, Markov's inequality and
Borel--Cantelli imply, along the dyadic sequence,
\[
 N_h\le h^{-1+(1/2-\varepsilon)\alpha-\zeta}
 \quad\text{eventually almost surely}.
\]
Each occupied cell contributes a set of diameter at most $\sqrt2r$ to
the $Z$-image and a set of diameter at most $r$ to a coordinate image.
Consequently their upper box dimensions are at most
\[
 \frac{1-(1/2-\varepsilon)\alpha+\zeta}
      {1/2-\varepsilon}.
\]
Letting $\varepsilon,\zeta$ decrease to zero through countable sequences
gives $2-\alpha$.  Taking $\alpha=2d_2$ or $\alpha=2$ proves the two
image assertions.  If $\alpha=2d_2+1>2$, choose $\varepsilon$ so small
that $(1/2-\varepsilon)\alpha>1$.  The right side of Equation
\eqref{arm:record-cell-count} is summable over dyadic $h$, so eventually
there are no candidate cells.  This proves the nonexistence assertion.
\end{proof}

\begin{lemma}[Generic horizontal contacts]\label{arm:generic-chords}
Almost surely, simultaneously on all compact time intervals:
\begin{enumerate}[label=\textup{(\roman*)}]
 \item The local minima of each coordinate are strict, countable, and
 have pairwise distinct values.  Every scalar horizontal equivalence
 class has at most three distinct times.
 \item No nontrivial horizontal chord is a chord of both $L$ and $R$.
 \item No member of a branch triple of one coordinate is an endpoint
 of a nontrivial horizontal chord of the other coordinate.
\end{enumerate}
The assertions also hold under any law having joint local absolute
continuity with these Brownian paths, including their relative coordinate
heights on separated windows.  In particular they hold in strict interior
windows of a nondegenerate correlated Brownian quadrant excursion.
\end{lemma}

\begin{proof}
We shall use the following precise source of the height densities.
Take deterministic intervals $I=[p,q]$ and $J=[r,v]$ with $q<r$.
The two paths
\[
 (Z(t)-Z(q))_{t\in I},\qquad (Z(t)-Z(r))_{t\in J}
\]
and the increment $G=Z(r)-Z(q)$ are independent.  The last has a
nonsingular Gaussian density on $\mathbb R^2$.  Its coordinate projections
have Gaussian densities on $\mathbb R$.  Thus, conditional on the two
displayed paths, a prescribed Lebesgue-null set of relative heights is
missed almost surely.  Conditioning is only on increments inside the
two separated windows; it is never on a putative contact or chord.

For completeness this also proves the scalar minimum facts.  The minimum
values on two disjoint compact rational intervals have probability zero
to agree: express the left minimum relative to its right endpoint and
the right minimum relative to its left endpoint, and condition on both
window paths.  Equality would prescribe one value for the independent
scalar Gaussian increment in the gap.  Taking a countable union shows
that no such pair of rational intervals has equal minima.  Two distinct
minimizers on a compact rational interval could be enclosed in disjoint
compact rational subintervals of that interval with the same minimum,
so each rational compact interval
has a unique minimizer.  Every local minimum is the minimizer on one
of these intervals.  This makes it strict, makes the set countable, and
also shows that distinct local minima have distinct values.  If a scalar
horizontal class contained four ordered visits, its middle two visits
would be distinct local minima at the same height.  This proves (i).

Suppose that $s<t$ were a chord in both coordinates.  Choose disjoint
rational windows $I,J$ containing $s,t$ in their interiors, and a
positive rational $\delta$ small enough that the forward arms from $s$
and the backward arms from $t$ of length $\delta$ lie in these windows.
Such choices are countable and cover every proposed nonzero chord.
Define compact sets of relative heights using all these local joint
record times in the respective windows:
\[
 A=\{Z(u)-Z(q):u\in I\text{ has both forward arms of length }\delta\},
\]
\[
 C=\{Z(w)-Z(r):w\in J\text{ has both backward arms of length }\delta\}.
\]
Here the record times are restricted to closed subintervals leaving
the indicated arm margins; countably many choices of these subintervals
suffice.  By Lemma~\ref{arm:record-height-images}, both sets have upper
box dimension at most $2(1-d_2)<1$.  Hence their difference set $A-C$
has two-dimensional Lebesgue measure zero.  Explicitly, for any
$a\in(2(1-d_2),1)$ each set has an $\epsilon$-cover with
$O(\epsilon^{-a})$ balls; pairwise differences cover $A-C$ with
$O(\epsilon^{-2a})$ balls, whose total area tends to zero.
The endpoint equality $Z(s)=Z(t)$ would imply $G\in A-C$.
The conditional Gaussian density therefore excludes the proposed
common chord, proving (ii).

Now let $a<b<c$ be an $L$-branch triple, and suppose an $R$-chord has
an endpoint $t\in\{a,b,c\}$ and another endpoint $u\ne t$.
If $t=b$, the $L$-coordinate has lower-wall arms on both sides of $b$,
and the $R$-chord supplies a lower-wall arm on the side pointing toward
$u$.  Truncate all three arms to a common positive rational length.
The final assertion of Lemma~\ref{arm:record-height-images} excludes
this case, since $d_2+1/2>1$.

If $t=a$ or $t=c$, retain a small window around $t$ disjoint from a
small window around $b$.  At $t$ the $L$-arm points toward $b$, and the
$R$-arm points toward $u$.  Truncate both to a common positive rational
length so that they lie in the outer-visit window.  They can point in
the same direction or in opposite directions; both possibilities are
included.  In relative $L$-heights, let $A$ be the image of all local
joint record times of this specified orientation in the outer-visit
window.  Lemma~\ref{arm:record-height-images} makes $A$ a
one-dimensional Lebesgue-null set: its upper box dimension is at most
$2(1-d_2)<1$ in the first case and zero in the second.

In the middle-visit window let $C$ be the countable set of minimum
values on its compact rational subintervals, measured relative to the
window's reference endpoint.  Since $b$ is a local minimum, its
relative $L$-height belongs to $C$.  Use the right endpoint as reference
for the earlier window and the left endpoint for the later window.
The equality $L(t)=L(b)$ would force the scalar Gaussian increment
between the windows to belong to either $A-C$ or $C-A$.
Both are Lebesgue-null, being countable unions of translates of a
null set.  Conditional on the two window paths the gap increment
has a density, so the equality has probability zero.  Countably many
rational windows, arm lengths, and orientations cover every outer
contact.  This proves (iii); interchanging the two coordinates gives
the other color.  The argument allows $u$ to equal another member of
the branch triple, as well as to lie between or outside its visits.

Every exclusion just proved is a null event specified in one window
or in finitely many separated windows together with their relative
heights.  It therefore transfers under the stated joint local absolute
continuity.  On strict interior windows a Brownian quadrant excursion
has this absolute continuity by its killed transition densities and
the Markov bridge construction, first restricting to compact positive
endpoint ranges and then exhausting them.  The excursion's common
initial/final endpoint is outside the assertion.
\end{proof}

\begin{remark}[Scope of the generic-contact assertion]
\label{arm:generic-chords-scope}
For an $L$-branch triple $a<b<c$, assertion (i) also gives
$L>L(b)$ on $(a,b)\cup(b,c)$ and lower values arbitrarily close just
outside $a$ and $c$.  Indeed an additional interior visit would give
four visits, while failure of an outside lower dip would make $a$ or
$c$ a second local minimum at the height of $b$.  Adjacent contacts
are consequently approached by ordinary chords at heights decreasing
to $L(b)$ from above, and the outer contact is approached by chords
at heights increasing to $L(b)$ from below.  This is a scalar consequence
of continuity and these strict inequalities.

The proof does not exclude a time that is incident to ordinary chords
of both colors with different other endpoints.  It excludes a common
pair and any other-color chord at a scalar branch triple.  In applying
the absolute-continuity clause, separate marginal absolute continuity
of the window increments is insufficient: the relative heights must
have the asserted joint law.  No statement is made at a singular center
where that joint local absolute continuity is unavailable.
\end{remark}


For later finite collections of viewed contour pieces, we need a version
which charges every possible choice of contact time.  The next argument
works directly with Brownian pieces and their relative heights.

\subsection{Polarity of finite Brownian chord diagrams}
\label{arm:finite-diagram-subsection}

Fix a positive definite covariance matrix $\Sigma$ whose two coordinate
variances agree and whose correlation is $\rho\in(0,1)$.  Put
\[
 d_0=0,\qquad d_1=\frac12,\qquad
 d_2=\frac{\pi}{2\arccos(-\rho)},\qquad
 \beta=\frac12+d_2>1.
\]
The argument below uses only Brownian survival and landing estimates.
It does not use a lattice arm estimate, a random address code, or a
passage kernel.

Let $J_i=[a_i^-,a_i^+]$, $1\le i\le m$, be pairwise disjoint compact
intervals in increasing order, and let $K_i$ be nonempty compact
subintervals of their interiors.  Fix $s>0$ smaller than half the
distance of every $K_i$ from $\partial J_i$.  On each $J_i$ take
\[
 X_i(t)=A_i+B_i(t-a_i^-),
\]
where the $B_i$ are independent Brownian motions with covariance
$\Sigma$.  The vector $(A_1,\ldots,A_m)$ may have any law independent
of these Brownian motions.  For absolute-continuity comparisons one
may take the $A_i$ to be independent nonsingular Gaussian vectors.

A \emph{completion} of these pieces is a continuous function
$f:[a_1^-,a_m^+]\to\mathbb R^2$ which agrees with $X_i$ on every
$J_i$.  No probability law for the missing gaps is imposed.
Let $G_1,G_2$ be graphs on $\{1,\ldots,m\}$, with loops discarded,
and write
\[
 e=\operatorname{rank}(G_1)+\operatorname{rank}(G_2),
 \qquad \operatorname{rank}(G_c)=m-\#\operatorname{components}(G_c),
\]
where isolated vertices count as components.  The graph is
\emph{realized} if there are $t_i\in K_i$ and a completion $f$ such
that, for each edge $\{i,j\}$ of $G_c$,
\begin{equation}
\label{arm:finite-diagram-chord}
 f_c(t_i)=f_c(t_j)=\min_{[t_i\wedge t_j,t_i\vee t_j]}f_c.
\end{equation}
Denote this event by $\mathcal E_G$.
In particular, a chord tests every retained part of its span, including
any intermediate retained windows.  Requiring only equality of its
endpoint heights would define a different event.

\begin{lemma}[Finite-cell Brownian diagram bound]
\label{arm:finite-diagram-bound}
\label{arm:finite-diagram-polarity}
For each $\varepsilon>0$ there are $C_\varepsilon<\infty$ and
$\delta_0>0$, depending on the fixed windows, $m$, $\Sigma$, and
$\varepsilon$, such that for $0<\delta<\delta_0$,
\begin{equation}
\label{arm:finite-diagram-probability}
 \mathbb P(\mathcal E_G)
 \le C_\varepsilon\delta^{\beta e-m-\varepsilon}.
\end{equation}
More precisely, there are measurable events
$\mathcal F_{G,\delta}\supset\mathcal E_G$ which satisfy the same
bound.  The constants are uniform in the law and values of the starting
vectors $A_i$.  Consequently $\mathcal E_G$ has probability zero
whenever $\beta e>m$.
\end{lemma}

\begin{proof}
We first state the elementary Brownian estimates used in the proof.
For a Brownian block of fixed duration $s$ and $k\in\{1,2\}$ constant
lower walls, let $\tau_k$ be the first violation, in the constrained
coordinates.  For $0<h\le1$,
\begin{align}
 \sup_{z\in[0,h]^k}\mathbb P_z(\tau_k>s)
       &\le C_s h^{2d_k},\label{arm:finite-diagram-survival}\\
 \sup_{z\in[0,\infty)^k}
 \mathbb P_z(\tau_k>s,\ Z(s)\in[0,h]^k)
       &\le C_s h^{k+2d_k}.\label{arm:finite-diagram-landing}
\end{align}
Fixed multiples of $h$ change only $C_s$.
Both bounds follow from Lemma~\ref{arm:killed-landing}, after translating
the constant walls to zero and fixing the positive duration.
The estimates remain valid conditionally when the walls and starting
point are known before the block is exposed.  An initial point below
a tested wall has zero probability of satisfying the gate.

Choose $\vartheta\in(0,1/2)$ with $2\beta e\vartheta<\varepsilon$, and put
$\eta=\delta^{1/2-\vartheta}$.  Cover each $K_i$ by intervals of a
deterministic mesh of width $\delta$, retaining the finitely many mesh
cells which meet $K_i$.  For sufficiently small $\delta$ these cells,
and an $s$-interval on either side of each cell, lie in $J_i$.
Let $\mathcal H_\delta$ be the event that both coordinate oscillations
on every such cell are at most $\eta$.  The reflection principle and
a union bound give
\begin{equation}
\label{arm:finite-diagram-mesh-error}
 \mathbb P(\mathcal H_\delta^c)
 \le C\delta^{-1}\exp\{-c\eta^2/\delta\}
 =C\delta^{-1}\exp\{-c\delta^{-2\vartheta}\}.
\end{equation}
This bound is uniform in all starting vectors.

Fix one cell $[a_i,b_i]$ for each $i$.  Suppose the graph is realized
with $t_i\in[a_i,b_i]$, and suppose $\mathcal H_\delta$ holds.
For one coordinate, the scalar minimum-chord relation is transitive.
Indeed all endpoints in a connected component have the same height,
and the union of the spans along a path between two vertices covers
the interval between their times; every such span stays above that
height.  Thus, in each component of $G_c$, the chronologically
consecutive selected times themselves form minimum chords.  Use
these sorted chains, with one comparison per unit of monochrome rank.
Let $k_i^-$ and $k_i^+$ be the numbers of colors with a predecessor
and successor, respectively, at vertex $i$.  Then
\begin{equation}
\label{arm:finite-diagram-incidences}
 k_i^-,k_i^+\in\{0,1,2\},\qquad
 \sum_i k_i^- =\sum_i k_i^+=e.
\end{equation}
Write $h_{i,c}=X_{i,c}(a_i)$.

Use the left slot $[a_i-s,a_i]$ and the right slot $[b_i,b_i+s]$.
All these slots are disjoint.  If color $c$ has predecessor $j$ at
$i$, let $H_c$ be their common actual chord height.  Then
$|h_{j,c}-H_c|\le\eta$ and $|h_{i,c}-H_c|\le\eta$.
The entire left slot lies between the two actual chord endpoints,
so, with the wall $w_{i,c}^-=h_{j,c}-\eta$,
\begin{equation}
\label{arm:finite-diagram-left-gate}
 \inf_{[a_i-s,a_i]}X_{i,c}\ge w_{i,c}^-,\qquad
 0\le X_{i,c}(a_i)-w_{i,c}^-\le3\eta.
\end{equation}
If color $c$ has a successor at $i$, the entire right slot likewise
lies in that chord span.  With $w_{i,c}^+=h_{i,c}-\eta$, cell
oscillation and the chord minimum give
\begin{equation}
\label{arm:finite-diagram-right-gate}
 \inf_{[b_i,b_i+s]}X_{i,c}\ge w_{i,c}^+,\qquad
 0\le X_{i,c}(b_i)-w_{i,c}^+\le2\eta.
\end{equation}
These are necessary gates involving only fixed-slot path restrictions,
not a conditioning on the actual selected times or on
$\mathcal H_\delta$.

Condition on all starting vectors and expose the independent Brownian
pieces in the order $J_1,\ldots,J_m$, each in increasing time.
Every left wall in Equation~\eqref{arm:finite-diagram-left-gate}
uses a representative in an earlier window.  Every right wall in
Equation~\eqref{arm:finite-diagram-right-gate} uses the already
exposed left endpoint of its cell.  Their heights are therefore
known before the respective slots.  Applying
Equations~\eqref{arm:finite-diagram-survival}
and~\eqref{arm:finite-diagram-landing} and the tower property bounds
the probability of all gates for this cell tuple by
\[
 C_m\eta^{\sum_i(k_i^-+2d_{k_i^-}+2d_{k_i^+})}.
\]
The bound is uniform in the starting vectors, so it also holds before
conditioning on them.  Since
$d_k\ge k d_2/2$ for $k\in\{0,1,2\}$,
Equation~\eqref{arm:finite-diagram-incidences} shows that the exponent
is at least
\[
 e+2d_2e=2\beta e.
\]
There are $O(\delta^{-m})$ tuples of cells.  Let
$\mathcal F_{G,\delta}$ be the union of the gate events for these
tuples, together with $\mathcal H_\delta^c$.  All its constituents
are measurable, and the preceding deterministic implications give
$\mathcal E_G\subset\mathcal F_{G,\delta}$.  Consequently
\[
 \mathbb P(\mathcal F_{G,\delta})
 \le C\delta^{\beta e-m-2\beta e\vartheta}
       +C\delta^{-1}e^{-c\delta^{-2\vartheta}}.
\]
Equation~\eqref{arm:finite-diagram-probability} follows by the choice
of $\vartheta$; the case $e=0$ also follows directly from the trivial bound
one.  If $\beta e>m$, choose $\varepsilon\in(0,\beta e-m)$ and
send $\delta$ to zero.

For completeness the realization event can also be taken measurable
without selecting a completion.  Replace every missing gap by affine
interpolation in each coordinate.  If any completion satisfies the
chords, every gap traversed by a chord has both its endpoint values
above that chord height, so affine interpolation also satisfies the
chords.  Thus a completion exists exactly when this fixed interpolation
has a realizing tuple.  Minimum values on compact intervals vary
continuously with the path and with the interval endpoints.  Taking
the minimum of the total nonnegative chord discrepancies over the
compact set $\prod_i K_i$ makes realization a measurable event.
\end{proof}

\begin{corollary}[Transfer and a finite class bound]
\label{arm:finite-diagram-transfer}
The polarity assertion of Lemma~\ref{arm:finite-diagram-bound} holds
under every law absolutely continuous with respect to the joint law
of the pieces and their starting vectors described above.
Under a density bound $L\le K$, its measurable mesh-cover events obey
\[
 \mathbb Q(\mathcal F_{G,\delta}\cap\{L\le K\})
 \le K C_\varepsilon\delta^{\beta e-m-\varepsilon}.
\]
Mere absolute continuity transfers polarity, without asserting the
same polynomial bound for the mesh-cover events under the untruncated
law.

Suppose a collection of eligible time points has a scalar chord
relation in each color, and every fixed finite family of pairwise
separated enlarged windows about distinct eligible points has this
joint absolute-continuity property.  Assume also that every tested
finite chord diagram is realized by a continuous completion in a
fixed order of those windows, or one of finitely many such orders.
Then, almost surely, every class generated by the two chord relations
has at most
\begin{equation}
\label{arm:finite-class-bound}
 \left\lfloor\frac{\beta}{\beta-1}\right\rfloor
\end{equation}
distinct eligible points, provided each finite connected subdiagram of
that class uses only eligible points and can be localized in a countable
family of such windows.
\end{corollary}

\begin{proof}
Absolute continuity transfers each null event.  The quantitative
truncated assertion is simply integration of the density over
$\mathcal F_{G,\delta}$.  There is no assertion of a uniform density
bound without the displayed truncation.

For a connected union $G_1\cup G_2$ on $m$ vertices, take a spanning
tree of the union.  Its edges of either color form an independent set
in that monochrome graph, so $e\ge m-1$.  If
$m>\beta/(\beta-1)$, then
$\beta e\ge\beta(m-1)>m$, and the diagram is polar.
There are finitely many two-colored graphs on a fixed finite vertex
set, and finitely many window orders.  Countable buffered localization
therefore excludes all such connected diagrams simultaneously.
If a generated class had at least
$m_0=\lfloor\beta/(\beta-1)\rfloor+1$ eligible vertices, finite chains
connecting them would give a finite connected graph on at least
$m_0$ vertices.  A spanning tree can be reduced by removing leaves to
a connected subgraph on exactly $m_0$ vertices.  Its edges are still
actual chords, contradicting the preceding exclusion.  This proves
Equation~\eqref{arm:finite-class-bound}.
\end{proof}

\begin{remark}[Separated windows on several lines]
\label{arm:finite-diagram-multiple-lines}
For finitely many off-center time points on finitely many lines, first
localize each point in a compact rational window bounded away from any
excluded center and from the other selected points on its line.  Retain
larger windows with positive buffers.  The argument applies whenever
their joint pair-path restrictions, including relative starting-height
vectors, are absolutely continuous with respect to independent
Brownian pieces with full-dimensional starting-height density.
An unrecorded common additive vector is harmless: the chord event is
translation invariant, and one can add an independent Gaussian common
vector when comparing full starting-height laws.  The chronological
order must be fixed in the finite diagram or enumerated among its
finitely many possibilities.

The completion requirement is substantive.  It includes lower-wall
conditions throughout every retained portion of a chord span; mere
endpoint equalities or endpoint-only arms do not supply these tests.
The assertion also requires joint absolute continuity of the pieces
and offsets, rather than just Brownian marginal laws or a density for
the offsets in isolation.  No claim is made at a pinned center lacking
this joint absolute continuity, or about chains through such a center.
\end{remark}

